\chapter{Communication Complexity — Trees}
%%% generated by the blueprint pipeline from TCSlib.CommunicationComplexity.DeterministicCC.Trees
\section{Overview}

This module develops the binary-tree structure underlying deterministic communication
protocols.  It introduces the \emph{shape} of a protocol as a \texttt{Tree Unit}, counts
the leaves of that tree, defines a subtree relation on binary trees, and proves the key
combinatorial fact that every binary tree with more than one leaf contains a balanced
subtree holding between one third and two thirds of all leaves.

%%%%%%%%%%%%%%%%%%%%%%%%%%%%%%%%%%%%%%%%%%%%%%%%%%%%%%%%%%%%%%%%%%%%%%%%%%%%%%%
\section{Declarations}

\begin{definition}[Protocol tree shape]
\lean{CommunicationComplexity.Deterministic.Protocol.shape}
\statementsource{roughgarden-cc-bootcamp}{
  pdf-d54d402b16e2-p004-b003,
  pdf-d54d402b16e2-p004-b005
}
\statementsource{raoyehudayoff-ch01-deterministic-protocols}{pdf-fa052d8c95c0-p003-b001}
\label{CommunicationComplexity.Deterministic.Protocol.shape}
\leanok
\uses{CommunicationComplexity.Deterministic.Protocol}
Given a deterministic communication protocol $p$ over input types $X$ and $Y$ with output
type $\alpha$, \texttt{shape} maps $p$ to its underlying binary tree of type
\texttt{Tree Unit}.  An output node becomes a leaf (\texttt{Tree.nil}), while an Alice or
Bob branching node becomes an internal node (\texttt{Tree.node}) whose two children are the
shapes of the two sub-protocols reached by the bit $\mathtt{false}$ and
$\mathtt{true}$, respectively.
\end{definition}

\begin{definition}[Number of protocol leaves]
\lean{CommunicationComplexity.Deterministic.Protocol.numLeaves}
\statementsource{roughgarden-cc-bootcamp}{pdf-d54d402b16e2-p004-b005}
\label{CommunicationComplexity.Deterministic.Protocol.numLeaves}
\leanok
\uses{CommunicationComplexity.Deterministic.Protocol, CommunicationComplexity.Deterministic.Protocol.shape}
The number of leaves (output nodes) of a protocol $p$ is defined as
$\texttt{numLeaves}(p) = \texttt{shape}(p).\texttt{numLeaves}$, i.e.\ the count of
\texttt{Tree.nil} nodes in the tree shape of $p$.
\end{definition}

\begin{definition}[Subtree relation on binary trees]
\lean{CommunicationComplexity.TreeIsSubtree}
\label{CommunicationComplexity.TreeIsSubtree}
\leanok
\texttt{TreeIsSubtree s t} is an inductively defined proposition asserting that $s$ is a
subtree of $t$.  It has two constructors: reflexivity ($s$ is a subtree of itself) and a
descent rule ($s$ is a subtree of the left child $\ell$ of $\texttt{node}(v,\ell,r)$,
hence a subtree of the whole node; an analogous right rule also exists in the
implementation).
\end{definition}

\begin{lemma}[Transitivity of the subtree relation]
\lean{CommunicationComplexity.TreeIsSubtree.trans}
\label{CommunicationComplexity.TreeIsSubtree.trans}
\leanok
\uses{CommunicationComplexity.TreeIsSubtree}
\difficulty{3}
Let $s$, $t$, and $u$ be binary trees with node labels in a type $\alpha$. If $s$ is a
subtree of $t$ and $t$ is a subtree of $u$, then $s$ is a subtree of $u$.
\end{lemma}

\begin{lemma}[Existence of a balanced subtree]
\lean{CommunicationComplexity.tree_balanced_subtree_aux}
\proofsource{raoyehudayoff-ch01-deterministic-protocols}{
  pdf-fa052d8c95c0-p004-b003,
  pdf-fa052d8c95c0-p004-b004
}
\label{CommunicationComplexity.tree_balanced_subtree_aux}
\leanok
\proofstep
  {CommunicationComplexity.Deterministic.Protocol}
  {direct}
  {raoyehudayoff-ch01-deterministic-protocols}
  {pdf-fa052d8c95c0-p003-b001}
\proofstep
  {CommunicationComplexity.Deterministic.Protocol.run}
  {direct}
  {raoyehudayoff-ch01-deterministic-protocols}
  {pdf-fa052d8c95c0-p003-b002}
\proofstep
  {CommunicationComplexity.Deterministic.Protocol.complexity}
  {direct}
  {raoyehudayoff-ch01-deterministic-protocols}
  {pdf-fa052d8c95c0-p003-b003}
\proofstep
  {CommunicationComplexity.Deterministic.Protocol.swap}
  {context}
  {raoyehudayoff-ch01-deterministic-protocols}
  {pdf-fa052d8c95c0-p003-b001}
\proofstep
  {CommunicationComplexity.Deterministic.Protocol.swap_run}
  {context}
  {raoyehudayoff-ch01-deterministic-protocols}
  {pdf-fa052d8c95c0-p003-b002}
\proofstep
  {CommunicationComplexity.Deterministic.Protocol.swap_complexity}
  {context}
  {raoyehudayoff-ch01-deterministic-protocols}
  {pdf-fa052d8c95c0-p003-b003}
\proofstep
  {CommunicationComplexity.Deterministic.Protocol.swap_swap}
  {context}
  {raoyehudayoff-ch01-deterministic-protocols}
  {pdf-fa052d8c95c0-p003-b001}
\proofstep
  {CommunicationComplexity.Deterministic.Protocol.comap}
  {context}
  {raoyehudayoff-ch01-deterministic-protocols}
  {pdf-fa052d8c95c0-p003-b001}
\proofstep
  {CommunicationComplexity.Deterministic.Protocol.comap_run}
  {context}
  {raoyehudayoff-ch01-deterministic-protocols}
  {pdf-fa052d8c95c0-p003-b002}
\proofstep
  {CommunicationComplexity.Deterministic.Protocol.comap_complexity}
  {context}
  {raoyehudayoff-ch01-deterministic-protocols}
  {pdf-fa052d8c95c0-p003-b003}
\proofstep
  {CommunicationComplexity.TreeIsSubtree}
  {context}
  {raoyehudayoff-ch01-deterministic-protocols}
  {pdf-fa052d8c95c0-p004-b003}
\proofstep
  {CommunicationComplexity.TreeIsSubtree.trans}
  {context}
  {raoyehudayoff-ch01-deterministic-protocols}
  {pdf-fa052d8c95c0-p004-b003}
\uses{CommunicationComplexity.TreeIsSubtree, CommunicationComplexity.TreeIsSubtree.trans}
\difficulty{5}
Let $t$ be a binary tree, and write $\ell(t)$ for its number of leaves. Let $n$ be a
natural number with $n > 1$, and suppose that $3\,\ell(t) \ge 2n$. Then $t$ has a
subtree $s$ whose number of leaves satisfies
\[
n \le 3\,\ell(s) < 2n.
\]
\end{lemma}

\begin{theorem}[Existence of a balanced subtree]
\lean{CommunicationComplexity.tree_balanced_subtree}
\proofsource{raoyehudayoff-ch01-deterministic-protocols}{
  pdf-fa052d8c95c0-p004-b003,
  pdf-fa052d8c95c0-p004-b004
}
\label{CommunicationComplexity.tree_balanced_subtree}
\leanok
\proofstep
  {CommunicationComplexity.Deterministic.Protocol}
  {direct}
  {raoyehudayoff-ch01-deterministic-protocols}
  {pdf-fa052d8c95c0-p003-b001}
\proofstep
  {CommunicationComplexity.Deterministic.Protocol.run}
  {direct}
  {raoyehudayoff-ch01-deterministic-protocols}
  {pdf-fa052d8c95c0-p003-b002}
\proofstep
  {CommunicationComplexity.Deterministic.Protocol.complexity}
  {direct}
  {raoyehudayoff-ch01-deterministic-protocols}
  {pdf-fa052d8c95c0-p003-b003}
\proofstep
  {CommunicationComplexity.Deterministic.Protocol.swap}
  {context}
  {raoyehudayoff-ch01-deterministic-protocols}
  {pdf-fa052d8c95c0-p003-b001}
\proofstep
  {CommunicationComplexity.Deterministic.Protocol.swap_run}
  {context}
  {raoyehudayoff-ch01-deterministic-protocols}
  {pdf-fa052d8c95c0-p003-b002}
\proofstep
  {CommunicationComplexity.Deterministic.Protocol.swap_complexity}
  {context}
  {raoyehudayoff-ch01-deterministic-protocols}
  {pdf-fa052d8c95c0-p003-b003}
\proofstep
  {CommunicationComplexity.Deterministic.Protocol.swap_swap}
  {context}
  {raoyehudayoff-ch01-deterministic-protocols}
  {pdf-fa052d8c95c0-p003-b001}
\proofstep
  {CommunicationComplexity.Deterministic.Protocol.comap}
  {context}
  {raoyehudayoff-ch01-deterministic-protocols}
  {pdf-fa052d8c95c0-p003-b001}
\proofstep
  {CommunicationComplexity.Deterministic.Protocol.comap_run}
  {context}
  {raoyehudayoff-ch01-deterministic-protocols}
  {pdf-fa052d8c95c0-p003-b002}
\proofstep
  {CommunicationComplexity.Deterministic.Protocol.comap_complexity}
  {context}
  {raoyehudayoff-ch01-deterministic-protocols}
  {pdf-fa052d8c95c0-p003-b003}
\proofstep
  {CommunicationComplexity.TreeIsSubtree}
  {context}
  {raoyehudayoff-ch01-deterministic-protocols}
  {pdf-fa052d8c95c0-p004-b003}
\proofstep
  {CommunicationComplexity.TreeIsSubtree.trans}
  {context}
  {raoyehudayoff-ch01-deterministic-protocols}
  {pdf-fa052d8c95c0-p004-b003}
\proofstep
  {CommunicationComplexity.tree_balanced_subtree_aux}
  {direct}
  {raoyehudayoff-ch01-deterministic-protocols}
  {pdf-fa052d8c95c0-p004-b003,
    pdf-fa052d8c95c0-p004-b004}
\uses{CommunicationComplexity.TreeIsSubtree, CommunicationComplexity.TreeIsSubtree.trans, CommunicationComplexity.tree_balanced_subtree_aux}
\difficulty{5}
Let $t$ be a binary tree with more than one leaf, and write $\abs{u}$ for the number of
leaves of a tree $u$. Then $t$ has a subtree $s$ whose leaf count satisfies
\[
  \abs{t} \;\le\; 3\abs{s} \;<\; 2\abs{t}.
\]
Equivalently, $s$ contains at least one third and strictly fewer than two thirds of the
leaves of $t$.
\end{theorem}
